% Targeted new A.22 attack, 2026-11-19; work draft, not a claimed solution.
\section*{An actual CM summand in a finite-cover Prym of the original punctured CM curve}
Let $K=\mathbf Q(i)$ and $E_0:y^2=x^3-x$. There is a degree-$32$ finite etale cover of $U=E_0\setminus\{O\}$ whose proper Jacobian has an $E_0$-isogenous factor in the kernel of pushforward to $E_0$. In particular, the residual weight-one alternative in LM6 A.22 cannot be removed by a general assertion that the proper Pryms of these covers have no CM characters.

Take the smooth plane quartic
\[
 C:\ x^4+y^4=z^4
\]
and the smooth proper genus-one curve $D$ with affine equation
\[
 s^2=1-t^4,\qquad O_D=(0,1).
\]
The rational formulas
\[
 q_x:C\longrightarrow D,\quad (t,s)=(y/z,x^2/z^2),\qquad
 q_y:C\longrightarrow D,\quad (t,s)=(x/z,y^2/z^2)
\]
extend uniquely to morphisms of smooth proper curves. Both have degree two: their function-field extensions adjoin respectively a square root of $s$, and the displayed sign-changing involutions are nontrivial deck involutions.

Here $D$ is the original CM elliptic curve over $K$, with a completely explicit isomorphism. Put
\[
 X=\frac{2(1+s)}{t^2},\qquad Y=\frac{4(1+s)}{t^3}.
\]
Then $Y^2=X^3+4X$. The inverse on the dense open set is
\[
 t=\frac{2X}{Y},\qquad s=\frac{2X^3}{Y^2}-1.
\]
Thus this is an isomorphism of smooth proper genus-one curves, sending $O_D$ to the point at infinity. With $a=1+i$, one has $a^4=-4$, so $X=a^2x_E$, $Y=a^3y_E$ identifies the cubic with $E_0$ over $K$.

The four branch values of $q_x$ are precisely $s=0$, $t^4=1$. Indeed the divisor of $s$ has simple zeroes at those four points and poles of order two at each of the two points at infinity of $D$; adjoining $\sqrt{s}$ therefore ramifies exactly at its zeroes. These branch values belong to $D[4]$. To see this without a torsion-coordinate assertion, let
\[
 j(t,s)=(-t,s),\qquad h(t,s)=(t,-s),\qquad P=(0,-1).
\]
The involution $j$ fixes the origin and is nontrivial, hence is $[-1]$ on the elliptic curve. It fixes $P$, so $2P=0$. The involution $h$ sends $O_D$ to $P$ and acts by $-1$ on the regular differential $dt/s$. An automorphism of an elliptic curve is the translation by its value at the origin followed by an origin-preserving automorphism; in characteristic zero the derivative determines the latter. Hence $h(Q)=P-Q$. Its fixed points, exactly the four branch values above, satisfy $2Q=P$ and thus $4Q=0$.

Consequently
\[
 f=[4]\circ q_x:C\longrightarrow D
\]
is a degree-$32$ morphism branched only over $O_D$, since $[4]$ is etale. It induces an actual finite etale map
\[
 C\setminus f^{-1}(O_D)\longrightarrow D\setminus\{O_D\}\simeq U.
\]
There are $28$ removed geometric points: four branch values of $q_x$ in $D[4]$ have one point above each, and the other twelve points of $D[4]$ have two. Thus $2-2g(C)-28=-32$, agreeing with the degree and $\chi(U)=-1$.

Write $q_{x,*}:J_C\to D$ for the norm map and $q_y^*:D\to J_C$ for pullback. The involution $\iota_x(x,y,z)=(-x,y,z)$ fixes $q_x$ and satisfies $q_y\iota_x=[-1]q_y$. Consequently
\[
 q_{x,*}\iota_x^*q_y^*=q_{x,*}q_y^*,\qquad
 \iota_x^*q_y^*=-q_y^*.
\]
It follows that $2q_{x,*}q_y^*=0$, hence
\[
 q_{x,*}q_y^*=0.
\]
The final implication is exact: the image of a morphism of abelian varieties is connected, whereas the image of a morphism annihilated by two lies in the finite etale group $D[2]$. Also $q_{y,*}q_y^*=[2]$, the norm--pullback identity for a degree-two map. Therefore, for every prime $\ell$,
\[
 q_y^*:T_\ell D\hookrightarrow T_\ell J_C,
 \qquad \operatorname{im}(q_y^*)\subset
 \ker\bigl(f_*:T_\ell J_C\to T_\ell D\bigr),
 \quad f_*=[4]q_{x,*}.
\]
Injection follows by applying $q_{y,*}$ and using torsion-freeness of $T_\ell D$, including $\ell=2$. For odd $\ell$, $\tfrac12q_{y,*}$ is an integral retraction. The connected image $q_y^*(D)$ lies in the connected Prym kernel of $q_{x,*}$ and is isogenous to $D\simeq E_0$.

All these maps are defined over $K$. For every rational prime $\ell$ split in $K$, the two idempotents of $K\otimes\mathbf Q_\ell\simeq\mathbf Q_\ell\times\mathbf Q_\ell$ give $G_K$-stable one-dimensional summands of $T_\ell D\otimes\mathbf Q_\ell$. Their images are actual weight-one CM lines in this proper Prym. This is a statement about a finite cover and its homology. It does not lift either line to a procyclic subgroup of $\pi_1(U_{\overline K})$, does not place such a lift in $R_{\rm nil}(F')$, and does not close A.22. The construction refutes precisely the suggested exclusion of all CM characters from all of the relevant finite-cover Pryms.

\subsection*{The actual punctured maps also eliminate this cover's proper homology alternative}
There is a stronger scoped consequence, using the pro-prime recognition theorem already proved in LM6 A.20 for $E_0$. The preceding CM Prym is not an example where one must artificially add punctures to obtain the elliptic quotient map: the required maps exist on the same open curve.

Let $Y=C\setminus f^{-1}(O_D)$. The map $q_y$ restricts to $Y\to D\setminus\{O_D\}$, since $q_y^{-1}(O_D)$ consists of $(0:1:1)$ and $(0:-1:1)$, whose $q_x$ images are the four-torsion points $(t,s)=(\pm1,0)$. The map $q_x$ itself also restricts to that punctured target because $q_x^{-1}(O_D)\subset f^{-1}(O_D)$.

The third quotient is
\[
 D_z:s_z^2=t_z^4+1,\quad O_z=(0,1),\qquad
 q_z:C\to D_z,\quad(t_z,s_z)=(x/y,z^2/y^2).
\]
Its origin has inverse image $(0:1:1),(0:1:-1)$, both in $f^{-1}(O_D)$. Thus $q_z$ restricts to $Y\to D_z\setminus\{O_z\}$ as well. The same formulas $X=2(1+s_z)/t_z^2$, $Y_E=4(1+s_z)/t_z^3$ give $Y_E^2=X^3-4X$. Scaling by $a=\sqrt2$, with $X=a^2x_E$, $Y_E=a^3y_E$, identifies $D_z$ with $E_0$ over $K'=K(\sqrt2)$. Passage to this finite field extension preserves the open arithmetic normalizer condition by taking the inverse image of $G_{K'}$.

The product of the three norm maps
\[
 J_C\longrightarrow D\times D\times D_z
\]
is an isogeny. Here is an explicit verification. On the affine chart $z=1$, write
\[
 \omega_1=\frac{dx}{y^3},\quad
 \omega_2=\frac{x\,dx}{y^3},\quad
 \omega_3=\frac{y\,dx}{y^3}.
\]
These are the usual three regular differentials on the smooth plane quartic. Differentiating $x^4+y^4=1$ gives $dy=-x^3dx/y^3$, and hence
\[
 q_x^*(dt/s)=-\omega_2,\qquad
 q_y^*(dt/s)=\omega_3,\qquad
 q_z^*(dt_z/s_z)=\omega_1.
\]
The pullback of differentials is an isomorphism. The product norm is consequently an isogeny in characteristic zero, and its map on the product of rational Tate modules is an isomorphism.

Let $A$ be a closed procyclic subgroup of the geometric fundamental group of $Y$, with open arithmetic normalizer image. For every prime $\ell$, its image in $H_1(C_{\overline K},\mathbf Z_\ell)$ is zero. Indeed, nonzero proper homology has nonzero image under one of the three quotient maps, by the isogeny just proved. The corresponding homomorphism of actual arithmetic fundamental groups carries $A$ to a procyclic subgroup of the maximal pro-$\ell$ geometric group of a copy of $E_0\setminus\{O\}$, and the normalizer still has open arithmetic image. The image is nonzero because its proper abelianization is nonzero. LM6 A.20 makes it cuspidal, whereas inertia at the unique puncture maps to zero in proper $H_1(E_0,\mathbf Z_\ell)$. This is a contradiction.

It follows in particular that whenever such an $A$ has nonzero image in $H_1(Y_{\overline K},\mathbf Z_\ell)$, that image lies in the boundary submodule. After a finite extension splitting the boundary points it is acted on through the cyclotomic character. A nonzero homomorphism from the torsion-free procyclic $A_\ell\simeq\mathbf Z_\ell$ to this torsion-free module is injective, so its action is the actual cyclotomic action on $A_\ell$. Nakamura's full profinite recognition theorem then places $A_\ell$ in a cusp inertia group of $Y$, and the cusp centralizer statement places all of $A$ there. This is an exact finite-cover homology recognition statement, not the assertion that this one cover detects every nontrivial subgroup in LM6 A.22.

\subsection*{Exact reach of the cover tower and a nonzero invisible subgroup}
Let $F=\pi_1(U_{\overline K})$, let $P_p=F^{(p)}$, and put
\[
 R=\bigcap_p\ker(F\to P_p),\qquad S=R_{\rm nil}(F').
\]
The subgroup $H=\pi_1(Y_{\overline K})\subset F$ contains $\ker(F\to P_2)$, and therefore contains every candidate $C\subset R$ in A.22. Indeed, on geometric function fields the degree-$16$ multiplication cover $[4]:D\to D$ is Galois with group $D[4]\simeq(\mathbf Z/4)^2$. Over its function field, $q_x$ is quadratic. The normal closure of the composition is obtained by adjoining the at most sixteen conjugate square roots. Its Galois group fits in
\[
 1\longrightarrow A\longrightarrow G\longrightarrow D[4]\longrightarrow1,
 \qquad A\subset(\mathbf Z/2)^{16},
\]
so $G$ is a finite $2$-group. The subgroup corresponding to the given connected cover contains the kernel of this quotient of $F$, and hence the kernel of the maximal pro-$2$ quotient.

The Fermat power maps give an explicit infinite family of further covers. For $N=4\cdot2^r$, put
\[
 C_N:\ x^N+y^N=z^N,\qquad
 a_N:C_N\to C,\quad(x:y:z)\mapsto(x^{N/4}:y^{N/4}:z^{N/4}),
 \qquad f_N=f\circ a_N.
\]
The map $a_N$ is Galois over $\overline K$ with group $(\mu_{N/4})^2$ and is ramified only above the coordinate-zero points of $C$. All those points map under $q_x$ to $D[4]$: $x=0$ was checked above; $y=0$ gives $O_D$ or $P\in D[2]$; and the two points of $D$ at infinity are fixed by $j=[-1]$, hence are in $D[2]$. Thus
\[
 Y_N=C_N\setminus f_N^{-1}(O_D)\longrightarrow U
\]
is finite etale. Iterating the normal-closure argument shows that its geometric Galois closure is a $2$-group. Every corresponding subgroup $H_N$ therefore contains $\ker(F\to P_2)$ and $R$.

This family cannot, as a purely geometric assertion, detect every nontrivial subgroup of the remaining kernel by ordinary first homology. In fact, for \emph{any} family of open subgroups $H_i\subset F$ containing $R$, one has the actual inclusions
\[
 R'\subset H_i',\qquad R'\longrightarrow H_i^{\rm ab}=0.
\]
Here primes mean closed commutator subgroups. The subgroup $S'$ is a nontrivial subgroup of $R'$ and already lies in the refined A.22 kernel $S$. To prove nontriviality, choose a surjection $\varphi:F\twoheadrightarrow A_5$, using a $5$-cycle and a $3$-cycle as images of the two free generators. Since $A_5$ is perfect, $\varphi(F')=A_5$. The quotient $F'/S$ is pronilpotent, so
\[
 A_5/\varphi(S)
\]
is a finite nilpotent quotient of the perfect group $A_5$, and is trivial. Hence $\varphi(S)=A_5$ and $\varphi(S')=A_5' = A_5$, proving $S'\ne1$. Moreover $S\subset R$: every map $F\to P_p$, restricted to $F'$, has pro-$p$ target, and therefore kills $S$.

Choose $g\in S'$ with $\varphi(g)\ne1$. Then $\overline{\langle g\rangle}$ is an actual nontrivial closed procyclic subgroup of $S'\subset S$, and its image in every $H_i^{\rm ab}$ is zero. This construction does \emph{not} assert that its arithmetic normalizer has open image. The missing arithmetic-normalizer requirement remains exactly the A.22 requirement. It proves the precise limitation of an argument using only the displayed covers, not a counterexample to A.22 itself. A successful finite-cover detector for all candidates must use the normalizer condition in an additional way, or range beyond this family.
